\subsubsection*{Antes de la sesión}
\begin{itemize}
\item[$\square$] Confirmar modalidad (Teams / híbrido / presencial) y materiales.
\item[$\square$] Revisar entregables de la sesión anterior (muestra).
\item[$\square$] Preparar demo en vivo con ejemplo STEM neutro EIC.
\item[$\square$] Recordar política de integridad / uso de IA (\ceddie).
\item[$\square$] Recordar presupuesto de trabajo independiente (bloques A--D; total 16\,h).
\end{itemize}

\subsubsection*{Mensajes clave}
\begin{enumerate}
\item \textbf{El brief es el contrato} entre docente y sistema.
\item \textbf{Orquestación $>$ prompt único} para obras largas o multi-artefacto.
\item \textbf{Criterio de parada} evita bucles y deja claro cuándo entra el humano.
\item Restricciones explícitas reducen ``ruido elegante''.
\end{enumerate}

\subsubsection*{Ruta sugerida (3\,h)}
\paragraph{Recap (10 min)} Conectar notas de integridad con el diseño del brief.
\paragraph{Brief (30 min)} Modelar brief pobre vs.\ bueno con ejemplo STEM neutro.
\paragraph{Prompt y roles (25 min)} Desglosar plantilla de prompt; introducir 3--4 roles típicos.
\paragraph{Demo (25 min)} Orquestación secuencial con dos aprobaciones humanas visibles.
\paragraph{Taller (50 min)} Priorizar calidad del brief y del diagrama sobre volumen de texto.
\paragraph{Cierre} Encargar TI bloque B; pedir idea de arquitectura de obra para sesión 3.

\subsubsection*{Contingencias}
\begin{itemize}
\item Si la herramienta falla: continuar con diseño en documento colaborativo.
\item Grupo heterogéneo: pista rápida (plantilla prellenada) y pista avanzada (multi-agente).
\item Dudas legales/políticas: anotar y derivar a CEDDIE; no improvisar normativa.
\end{itemize}

\subsubsection*{Checklist post-sesión}
\begin{itemize}
\item[$\square$] Recordar canal de entrega del producto de la sesión.
\item[$\square$] Asignar / recordar TI del intervalo (ver cronograma parte B).
\item[$\square$] Anotar fricciones de la demo para la siguiente sesión.
\end{itemize}
